\documentclass[twocolumn,10pt,a4paper]{article}

% --- Pacotes Fundamentais ---
\usepackage[utf8]{inputenc}
\usepackage[T1]{fontenc}
\usepackage[brazil]{babel}

% --- Tipografia (Sans-Serif Idêntica ao Modelo da Kamily) ---
\usepackage{helvet}
\renewcommand{\familydefault}{\sfdefault}

% --- Margens e Duas Colunas ---
\usepackage[top=2cm, bottom=2cm, left=1.8cm, right=1.8cm, columnsep=0.75cm]{geometry}

% --- Pacotes de Matemática, Tabelas e Figuras ---
\usepackage{amsmath, amssymb}
\usepackage{booktabs}
\usepackage{graphicx}
\graphicspath{{figures/}}

% --- Formatação de Seções e Legendas ---
\usepackage{caption}
\captionsetup{font=footnotesize, labelfont=bf, skip=4pt}

\usepackage{titlesec}
\titleformat{\section}{\large\bfseries}{}{0pt}{}
\titleformat{\subsection}{\normalsize\bfseries}{}{0pt}{}
\titlespacing*{\section}{0pt}{1.2ex plus 0.4ex minus 0.2ex}{0.5ex plus 0.2ex}
\titlespacing*{\subsection}{0pt}{1ex plus 0.3ex minus 0.2ex}{0.4ex plus 0.2ex}

% Espaçamento entre parágrafos
\setlength{\parindent}{0pt}
\setlength{\parskip}{0.55em}

\begin{document}

% ==============================================================================
% RESUMO (Início direto sem título nem cabeçalho, exatamente como no modelo)
% ==============================================================================
\section*{Resumo}

O Exame Nacional do Ensino Médio (ENEM) é a principal porta de entrada para o ensino superior no Brasil, e seus resultados costumam refletir desigualdades socioeconômicas e educacionais. Este trabalho investiga como fatores familiares (escolaridade dos pais) e escolares (rede, localização e infraestrutura), além do contexto econômico do DF, se associam ao desempenho em Matemática dos estudantes do Distrito Federal no ENEM entre 2023 e 2025, e em que medida essas relações variam entre grupos e contextos. Utilizam-se os microdados do ENEM, o Censo Escolar da Educação Básica e os Indicadores Educacionais das Escolas (INSE, Adequação Docente e Esforço Docente), do Inep, além do PIB dos Municípios, do IBGE. A análise combina exploração visual (mapas de calor e gráficos de barras agrupadas), regressão linear múltipla e Random Forest Regressor, tendo a nota de Matemática como variável dependente. Espera-se identificar quais fatores têm maior peso relativo no desempenho e como esse peso se altera conforme o perfil do estudante e da escola, de modo a oferecer subsídios para discussões sobre equidade educacional.

\textbf{Palavras-chave:} ENEM; Matemática; desigualdade educacional; nível socioeconômico; mineração de dados educacionais; Random Forest.

% ==============================================================================
% 1. INTRODUÇÃO
% ==============================================================================
\section*{1. Introdução}

O ENEM, criado em 1998, tornou-se o principal mecanismo de acesso ao ensino superior brasileiro, sendo usado por universidades públicas via SISU e por programas como o ProUni e o Fies. Por reunir milhões de participantes todos os anos, também é uma das fontes mais ricas para estudar o desempenho de estudantes ao fim da educação básica e as desigualdades que os cercam.

A literatura mostra que o desempenho no exame não depende só do esforço individual. Estudos com microdados do ENEM apontam que o desempenho dos estudantes frequentemente se correlaciona com o background familiar, o tipo de escola frequentada e o acesso a recursos educacionais, e há trabalhos recentes que aplicam aprendizado de máquina a esses dados em recortes regionais. Dentro dessa discussão, a escolaridade dos pais é recorrentemente apontada como um dos principais preditores do sucesso acadêmico.

No cenário brasileiro, a diferença entre escolas públicas e privadas é um dos eixos mais discutidos. Na avaliação PISA 2018, escolas privadas brasileiras tiveram resultados bem melhores que as públicas em Matemática, Ciências e Leitura (OECD, 2018), e Moraes e Belluzzo (2014) também documentam esse diferencial. Esse debate motivou políticas de ação afirmativa, como a Lei de Cotas (Lei nº 12.711/2012), que ampliou o acesso de egressos da escola pública às universidades federais, mas atua apenas no ingresso. Pesquisas sobre o ensino superior mostram que as desvantagens de preparo e de condições socioeconômicas continuam a afetar o desempenho depois da entrada, o que reforça a importância de entender essas desigualdades já na educação básica.

O Distrito Federal é um recorte especialmente interessante. Reúne uma das maiores rendas per capita do país e, ao mesmo tempo, forte desigualdade entre suas regiões administrativas, com redes pública e privada atuando no mesmo território. Isso permite estudar desigualdade educacional em um contexto único, sem a heterogeneidade regional de análises nacionais. Além disso, poucos estudos com microdados do ENEM tratam especificamente do DF.

Optou-se por analisar exclusivamente o desempenho em Matemática, em razão do vínculo dessa área com a educação em computação. Como o ENEM é a principal via de ingresso nos cursos de graduação, o que inclui áreas da tecnologia, entender quais fatores familiares, escolares e territoriais se associam à nota de Matemática ajuda a identificar desigualdades de preparo que chegam à graduação. Além disso, concentrar a análise em uma única nota evita misturar fatores que afetam cada área de modo diferente e deixa as comparações entre grupos mais claras.

Assim, a questão de pesquisa é: \textbf{Como fatores socioeconômicos e educacionais, incluindo a escolaridade dos pais e as características da escola (rede, localização e infraestrutura), se associam ao desempenho em Matemática dos estudantes do Distrito Federal no ENEM, e em que medida essas relações variam entre diferentes grupos e contextos (2023--2025)?}

Para respondê-la, usamos dados abertos do Inep e do IBGE e três etapas: exploração visual (mapas de calor e barras agrupadas), regressão linear múltipla e Random Forest Regressor para ranquear a importância relativa dos fatores.

% ==============================================================================
% 2. BACKGROUND
% ==============================================================================
\section*{2. Background}

\subsection*{2.1 Origem social, escola e desigualdade de desempenho}

A ideia de que o desempenho escolar reflete a origem social tem longa tradição na sociologia da educação. Bourdieu e Passeron (1970) veem a escola como um espaço que tende a reproduzir desigualdades, porque o desempenho é fortemente influenciado pelo capital cultural das famílias. Nessa linha, a escolaridade dos pais funciona como indicador de capital cultural e de apoio educacional em casa, e é, junto com a renda, uma das variáveis mais usadas em estudos de desempenho no ENEM. Análises com regressão para o Ceará encontraram desempenho maior entre participantes com pais de maior escolaridade. Trabalhos relacionados com o Nordeste também indicam que parte relevante da variação das notas é explicada pelas diferenças entre escolas, e que escolaridade dos pais e renda têm impacto positivo no desempenho.

\subsection*{2.2 Rede, localização e infraestrutura escolar}

Além da família, características da escola importam. A diferença entre redes pública e privada aparece em avaliações nacionais e internacionais, mas parte dela vem do perfil socioeconômico dos alunos que cada rede atende, e não só da escola em si. Por isso é comum usar indicadores como o INSE, que resume o nível socioeconômico das escolas, ao lado de variáveis como localização (urbana ou rural) e infraestrutura (laboratórios, acesso à internet), além de indicadores docentes como Adequação e Esforço Docente.

\subsection*{2.3 Contexto territorial e econômico}

O território onde a escola está também condiciona oportunidades: renda, acesso a serviços e infraestrutura variam entre regiões. No DF, essa variação aparece principalmente entre as regiões administrativas, que diferem bastante em renda e perfil socioeconômico. Estudos sobre o ENEM mostram diferenças regionais marcantes em outros recortes. Incluir um indicador de renda por região administrativa permite testar se o contexto territorial se associa às notas além do que o INSE e as características da escola já explicam.

\subsection*{2.4 Mineração de dados educacionais e modelos preditivos}

A Mineração de Dados Educacionais (EDM) aplica técnicas de mineração de dados a bases educacionais para gerar conhecimento que apoie políticas e práticas de ensino. Entre as técnicas, a regressão linear múltipla é útil por ser interpretável: seus coeficientes indicam o efeito de cada fator, mantendo os demais constantes. Já o Random Forest captura relações não lineares e interações entre variáveis e permite medir a importância relativa de cada preditor, que se aproxima da nossa proposta de comparar contextos. No ENEM, há aplicações recentes de aprendizado de máquina com técnicas de explicabilidade para identificar os fatores mais relevantes.

Neste trabalho, os dois métodos se complementam: a regressão mede o efeito simultâneo dos fatores sobre a nota de Matemática, e o Random Forest ranqueia a importância relativa deles, dando duas visões sobre a mesma pergunta.

\subsection*{2.5 Trabalhos Relacionados}

No âmbito da educação em computação no Distrito Federal, destaca-se a pesquisa seminal de Silveira, Holanda e Ramos (2022), intitulada \textit{"Analysis of Student Performance and Social-economic Data in Introductory Computer Science Courses at the University of Brasília"}, apresentada na conferência IEEE Frontiers in Education (FIE). Os autores aplicaram técnicas de Mineração de Dados Educacionais e Modelos Lineares Generalizados (GLM) aos dados de calouros matriculados na disciplina de introdução à programação (CS1) nos quatro cursos do Departamento de Ciência da Computação da UnB entre 2012 e 2019.

Aquele estudo revelou achados cruciais para o entendimento da evasão e da retenção acadêmica na UnB: o egresso de escola privada apresentou coeficiente positivo e significante no desempenho em programação (+0,796), ao passo que a vulnerabilidade socioeconômica associou-se negativamente à aprovação (-0,768). O presente trabalho dialoga diretamente com esses achados, recuando uma etapa estratégica: investigamos a formação e os filtros na conclusão do ensino médio no próprio DF, identificando a magnitude com que o capital cultural, o tipo de escola e a infraestrutura tecnológica domiciliar estabelecem os diferenciais que chegam à graduação.

% ==============================================================================
% 3.0 METODOLOGIA
% ==============================================================================
\section*{3.0 Metodologia}

\subsection*{3.1 Amostra e Coleta}

A base de dados integra três fontes abertas oficiais referentes ao triênio 2021--2023: os microdados do ENEM (Inep), os microdados do Censo Escolar da Educação Básica (Inep) e o Produto Interno Bruto dos Municípios (IBGE). Para viabilizar a manipulação dos arquivos volumosos, implementou-se um pipeline em Python que realiza a leitura e filtragem em fluxo contínuo (streaming por blocos de 100.000 linhas) diretamente a partir dos arquivos compactados (.zip), sem descompressão integral em disco.

Foram aplicados três critérios estritos de corte e filtragem amostral:
\begin{enumerate}
  \item Recorte geográfico exclusivo do Distrito Federal (\texttt{SG\_UF\_PROVA == 'DF'});
  \item Presença efetiva no segundo dia de prova com nota válida em Matemática (\texttt{TP\_PRESENCA\_MT == 1});
  \item Preenchimento das variáveis socioeconômicas e familiares indispensáveis.
\end{enumerate}
A coorte final consolidada para modelagem resultou em $N = 126.874$ estudantes do Distrito Federal.

\subsection*{3.2 Variáveis e Engenharia de Atributos}

A variável dependente ($Y$) é a Nota de Matemática e suas Tecnologias (\texttt{NU\_NOTA\_MT}). Os preditores ($X$) foram submetidos à engenharia de atributos:
\begin{itemize}
  \item A escolaridade do pai (\texttt{Q001}) e da mãe (\texttt{Q002}) foi convertida em anos de escolarização formais (0 a 18 anos), extraindo-se a Maior Escolaridade dos Pais (\texttt{MAX\_ESC\_PAIS});
  \item A renda familiar (\texttt{Q006}) foi convertida para escala logarítmica ($\log(\text{Renda})$) para mitigar a assimetria da distribuição;
  \item Foram criadas variáveis binárias para posse de computador (\texttt{Q024}) e internet domiciliar (\texttt{Q025});
  \item A dependência administrativa foi codificada em variáveis dummy (Escola Privada, Escola Federal e Escola Pública Distrital);
  \item Variáveis de controle incluíram pertença étnico-racial (PPI), sexo feminino e localização urbana/rural.
\end{itemize}

\subsection*{3.3 Modelagem Estatística e Preditiva}

A estratégia analítica combina duas perspectivas complementares: Regressão Linear Múltipla estimada por Mínimos Quadrados Ordinários (OLS), para inferência paramétrica e medição de efeitos marginais \textit{ceteris paribus}; e Random Forest Regressor supervisionado com 100 árvores de decisão e profundidade controlada ($\text{max\_depth} = 15$). O conjunto amostral foi particionado em 70\% para treinamento ($N = 88.811$) e 30\% para teste ($N = 38.063$), com $\text{random\_state} = 42$. A importância relativa de cada atributo foi computada via \textit{Mean Decrease in Impurity} (MDI).

% ==============================================================================
% 4.0 RESULTADOS E DISCUSSÃO
% ==============================================================================
\section*{4.0 Resultados e Discussão}

\subsection*{4.1 Análise Exploratória e Visual}

A Figura~1 apresenta a matriz de correlação de Pearson entre os fatores investigados e o rendimento em Matemática. A nota exibe forte correlação positiva com a Renda Familiar Mensal ($r = 0{,}47$), a Escolaridade do Pai ($r = 0{,}38$), a Escolaridade da Mãe ($r = 0{,}36$) e a Posse de Computador ($r = 0{,}33$), contrastando com correlação negativa com a condição PPI ($r = -0{,}26$).

\begin{figure}[htbp]
  \centering
  \includegraphics[width=\columnwidth]{figures/heatmap_correlacoes.png}
  \caption{Matriz de Correlação de Pearson: Variáveis Inep, IBGE e Nota de Matemática.}
\end{figure}

Conforme previsto na formulação do projeto, a Figura~2 apresenta a mediana das notas em Matemática isolando as redes de ensino em cada nível de estudo dos genitores. Observa-se que, enquanto alunos da rede pública de pais sem instrução obtêm mediana de 435,5 pontos, alunos de pais com pós-graduação atingem 565 pontos. Contudo, em qualquer patamar de escolaridade familiar, a escola privada confere um salto marcante, atingindo mediana superior a 630 pontos nos estratos mais altos (mediana geral de 629,9 pontos versus 501,3 da rede distrital).

\begin{figure}[htbp]
  \centering
  \includegraphics[width=\columnwidth]{figures/barras_mediana_escolaridade_rede.png}
  \caption{Mediana em Matemática por Maior Escolaridade dos Pais e Rede de Ensino no DF.}
\end{figure}

A Figura~3 detalha a mediana das notas conforme a disponibilidade de computador no domicílio. Os participantes com computador alcançam mediana agregada de 575,5 pontos, contra 474,3 pontos dos que não dispõem do equipamento (diferencial expressivo de 101,2 pontos). Esse padrão ratifica que os bens tecnológicos domésticos ampliam a capacidade de fixação de conteúdos quantitativos e autoinstrução.

\begin{figure}[htbp]
  \centering
  \includegraphics[width=\columnwidth]{figures/barras_mediana_escolaridade_computador.png}
  \caption{Mediana em Matemática por Escolaridade dos Pais e Posse de Computador no DF.}
\end{figure}

\subsection*{4.2 Efeitos Marginais: Regressão Linear Múltipla}

A Tabela~1 reúne os coeficientes da regressão linear múltipla (OLS). O modelo alcançou $R^2 = 0{,}2784$ com elevada significância ($F = 4.079, p < 0{,}001$), demonstrando que mais de um quarto da variabilidade das notas em Matemática decorre exclusivamente das variáveis estruturais e familiares inseridas.

\begin{table}[htbp]
  \centering
  \footnotesize
  \caption{Coeficientes da Regressão Linear Múltipla OLS ($R^2 = 0{,}2784$, $N = 126.874$).}
  \begin{tabular}{l rrrr}
    \toprule
    \textbf{Variável} & \textbf{Coef. ($\beta$)} & \textbf{Erro P.} & \textbf{$t$} & \textbf{$p$-valor} \\
    \midrule
    Intercepto ($\beta_0$)   & $373{,}13$ & $2{,}17$ & $171{,}6$ & $< 0{,}001$ \\
    Esc. Pai (anos)          & $+4{,}08$  & $0{,}09$ & $46{,}2$  & $< 0{,}001$ \\
    Esc. Mãe (anos)          & $+3{,}65$  & $0{,}09$ & $38{,}9$  & $< 0{,}001$ \\
    $\log(\text{Renda})$     & $+13{,}51$ & $0{,}22$ & $61{,}7$  & $< 0{,}001$ \\
    Tem Computador           & $+27{,}46$ & $0{,}84$ & $32{,}7$  & $< 0{,}001$ \\
    Escola Privada           & $+44{,}20$ & $1{,}15$ & $38{,}3$  & $< 0{,}001$ \\
    Escola Federal           & $+45{,}39$ & $2{,}53$ & $17{,}9$  & $< 0{,}001$ \\
    Cor/Raça PPI             & $-24{,}19$ & $0{,}65$ & $-37{,}1$ & $< 0{,}001$ \\
    Sexo Feminino            & $-32{,}85$ & $0{,}64$ & $-51{,}5$ & $< 0{,}001$ \\
    \bottomrule
  \end{tabular}
\end{table}

Os coeficientes revelam que cada ano de estudo adicional do pai e da mãe agrega, respectivamente, +4,08 e +3,65 pontos líquidos à nota. Destaca-se o prêmio institucional da escola privada (+44,20 pontos) e federal (+45,39 pontos) frente à pública distrital, comprovando que o diferencial de desempenho persiste mesmo isolando a renda domiciliar.

\subsection*{4.3 Preditores Não Lineares e Ranking: Random Forest}

No conjunto de teste (30\%), o Random Forest obteve $R^2 = 0{,}2896$, com erro absoluto médio (MAE) de 85,60 pontos. A Figura~4 expõe o ranking de importância de variáveis (MDI): a Renda Familiar Mensal lidera com 64,05\% da força preditiva, seguida pela escolaridade paterna (8,22\%) e materna (7,79\%), somando 16,01\% para o capital cultural. Dispositivos digitais, sexo e tipologia escolar complementam a capacidade preditiva.

\begin{figure}[htbp]
  \centering
  \includegraphics[width=\columnwidth]{figures/ranking_importancia_variaveis.png}
  \caption{Ranking de Importância de Variáveis via Random Forest Regressor (MDI).}
\end{figure}

A Tabela~2 detalha o comportamento médio e mediano das notas de Matemática desagregadas pelos estratos sociodemográficos e escolares investigados no Distrito Federal.

\begin{table}[htbp]
  \centering
  \footnotesize
  \caption{Estatísticas Descritivas de Matemática por Subgrupo no DF ($N = 126.874$).}
  \begin{tabular}{l rrr}
    \toprule
    \textbf{Dimensão / Subcategoria} & \textbf{$N$} & \textbf{Média} & \textbf{Mediana} \\
    \midrule
    \multicolumn{4}{l}{\textit{\textbf{Rede de Ensino}}} \\
    \quad Rede Privada             & $18.377$  & $621{,}1$ & $629{,}9$ \\
    \quad Rede Federal (IFB)       & $2.207$   & $601{,}0$ & $600{,}5$ \\
    \quad Rede Pública Distrital   & $26.063$  & $511{,}6$ & $501{,}3$ \\
    \addlinespace
    \multicolumn{4}{l}{\textit{\textbf{Escolaridade dos Pais}}} \\
    \quad Pais: Pós-graduação      & $39.018$  & $623{,}7$ & $633{,}9$ \\
    \quad Pais: Ensino Superior    & $27.731$  & $586{,}2$ & $588{,}1$ \\
    \quad Pais: Ensino Médio       & $42.617$  & $521{,}1$ & $509{,}5$ \\
    \quad Pais: Sem estudo         & $1.574$   & $448{,}8$ & $435{,}5$ \\
    \addlinespace
    \multicolumn{4}{l}{\textit{\textbf{Recursos Tecnológicos}}} \\
    \quad Possui Computador        & $105.100$ & $578{,}0$ & $575{,}5$ \\
    \quad Não Possui Computador    & $35.349$  & $487{,}5$ & $474{,}3$ \\
    \bottomrule
  \end{tabular}
\end{table}

\end{document}
